\documentclass[exercice]{thedocument}%
\startdatakeys
\source{}
\BO{2026}
\cycle{4}
\competencesDuSocle{RE}
\objectifsApprentissage{B1ca,B1cb}
\automatismes{B1ca,B2da,B2db}
\enddatakeys
\begin{document}%
\begin{enumerate}
    \item Donner la définition de \textbf{deux figures symétriques par
    rapport à une droite}.
    \dotedline\dotedline\dotedline\dotedline
    \item Si A et A' sont deux points symétriques par rapport à une droite
    (d), donner \textbf{deux propriétés} liant la droite (d) et le segment [AA'].
    \dotedline\dotedline\dotedline\dotedline
    \item Écrire la méthode permettant de construire le symétrique A' d'un
    point A par rapport à une droite (d).
    \dotedline\dotedline\dotedline\dotedline
    \item Donner deux grandeurs qui sont conservées par symétrie axiale.
    \dotedline\dotedline\dotedline\dotedline
\end{enumerate}
\end{document}
